\IfFileExists{stacks-project.cls}{%
\documentclass{stacks-project}
}{%
\documentclass{amsart}
}
\usepackage{amssymb}

% For dealing with references we use the comment environment
\usepackage{verbatim}
\usepackage{fancyvrb}
\usepackage{pmboxdraw}
\newenvironment{reference}{\comment}{\endcomment}
%\newenvironment{reference}{}{}
\newenvironment{slogan}{\comment}{\endcomment}
\newenvironment{history}{\comment}{\endcomment}

% For source figures retained as assets
\usepackage{graphicx}

% For commutative diagrams we use Xy-pic
\usepackage[all]{xy}

% We use 2cell for 2-commutative diagrams.
\xyoption{2cell}
\UseAllTwocells

% We use multicol for the list of chapters between chapters
\usepackage{multicol}

% This is generally recommended for better output
\usepackage{lmodern}
\usepackage[T1]{fontenc}
\usepackage{textalpha}

% For cross-file-references
% We prefer the snapshot tmp/xref/#2.aux,
% which latexmkrc replaces atomically, so
% that documents compiled simultaneously
% never read a half-written .aux file.
% The last argument keeps links to #2.pdf.
\usepackage{xr-hyper}
\newcommand{\maybeexternaldocument}[2]{%
  \IfFileExists{#2.aux}{%
    \externaldocument[#1]{#2}%
  }{%
    \PackageWarning{stack}{Missing #2.aux; external references with prefix #1 unavailable}%
  }%
}
\let\externaldocumentorig\externaldocument
\renewcommand{\externaldocument}[2][]{%
  \IfFileExists{tmp/xref/#2.aux}{%
    \externaldocumentorig[#1]{tmp/xref/#2}[#2.pdf]%
  }{%
  \IfFileExists{#2.aux}{%
    \externaldocumentorig[#1]{#2}%
  }{%
    \PackageWarning{stack}{Missing #2.aux; external references with prefix #1 unavailable}%
  }%
  }%
}

% Package for hypertext links:
\usepackage{hyperref}

% For any local file, say "hello.tex" you want to link to please
% use \externaldocument[hello-]{hello}
\externaldocument[introduction-]{introduction}
\externaldocument[conventions-]{conventions}
\externaldocument[sets-]{sets}
\externaldocument[categories-]{categories}
\externaldocument[topology-]{topology}
\externaldocument[sheaves-]{sheaves}
\externaldocument[sites-]{sites}
\externaldocument[stacks-]{stacks}
\externaldocument[fields-]{fields}
\externaldocument[algebra-]{algebra}
\externaldocument[brauer-]{brauer}
\externaldocument[homology-]{homology}
\externaldocument[derived-]{derived}
\externaldocument[simplicial-]{simplicial}
\externaldocument[more-algebra-]{more-algebra}
\externaldocument[smoothing-]{smoothing}
\externaldocument[modules-]{modules}
\externaldocument[sites-modules-]{sites-modules}
\externaldocument[injectives-]{injectives}
\externaldocument[cohomology-]{cohomology}
\externaldocument[sites-cohomology-]{sites-cohomology}
\externaldocument[dga-]{dga}
\externaldocument[dpa-]{dpa}
\externaldocument[sdga-]{sdga}
\externaldocument[hypercovering-]{hypercovering}
\externaldocument[schemes-]{schemes}
\externaldocument[constructions-]{constructions}
\externaldocument[properties-]{properties}
\externaldocument[morphisms-]{morphisms}
\externaldocument[coherent-]{coherent}
\externaldocument[divisors-]{divisors}
\externaldocument[limits-]{limits}
\externaldocument[varieties-]{varieties}
\externaldocument[topologies-]{topologies}
\externaldocument[descent-]{descent}
\externaldocument[perfect-]{perfect}
\externaldocument[more-morphisms-]{more-morphisms}
\externaldocument[flat-]{flat}
\externaldocument[groupoids-]{groupoids}
\externaldocument[more-groupoids-]{more-groupoids}
\externaldocument[etale-]{etale}
\externaldocument[chow-]{chow}
\externaldocument[intersection-]{intersection}
\externaldocument[pic-]{pic}
\externaldocument[weil-]{weil}
\externaldocument[adequate-]{adequate}
\externaldocument[dualizing-]{dualizing}
\externaldocument[duality-]{duality}
\externaldocument[discriminant-]{discriminant}
\externaldocument[derham-]{derham}
\externaldocument[local-cohomology-]{local-cohomology}
\externaldocument[algebraization-]{algebraization}
\externaldocument[curves-]{curves}
\externaldocument[resolve-]{resolve}
\externaldocument[models-]{models}
\externaldocument[functors-]{functors}
\externaldocument[equiv-]{equiv}
\externaldocument[pione-]{pione}
\externaldocument[etale-cohomology-]{etale-cohomology}
\externaldocument[proetale-]{proetale}
\externaldocument[relative-cycles-]{relative-cycles}
\externaldocument[more-etale-]{more-etale}
\externaldocument[trace-]{trace}
\externaldocument[crystalline-]{crystalline}
\externaldocument[spaces-]{spaces}
\externaldocument[spaces-properties-]{spaces-properties}
\externaldocument[spaces-morphisms-]{spaces-morphisms}
\externaldocument[decent-spaces-]{decent-spaces}
\externaldocument[spaces-cohomology-]{spaces-cohomology}
\externaldocument[spaces-limits-]{spaces-limits}
\externaldocument[spaces-divisors-]{spaces-divisors}
\externaldocument[spaces-over-fields-]{spaces-over-fields}
\externaldocument[spaces-topologies-]{spaces-topologies}
\externaldocument[spaces-descent-]{spaces-descent}
\externaldocument[spaces-perfect-]{spaces-perfect}
\externaldocument[spaces-more-morphisms-]{spaces-more-morphisms}
\externaldocument[spaces-flat-]{spaces-flat}
\externaldocument[spaces-groupoids-]{spaces-groupoids}
\externaldocument[spaces-more-groupoids-]{spaces-more-groupoids}
\externaldocument[bootstrap-]{bootstrap}
\externaldocument[spaces-pushouts-]{spaces-pushouts}
\externaldocument[spaces-chow-]{spaces-chow}
\externaldocument[groupoids-quotients-]{groupoids-quotients}
\externaldocument[spaces-more-cohomology-]{spaces-more-cohomology}
\externaldocument[spaces-simplicial-]{spaces-simplicial}
\externaldocument[spaces-duality-]{spaces-duality}
\externaldocument[formal-spaces-]{formal-spaces}
\externaldocument[restricted-]{restricted}
\externaldocument[spaces-resolve-]{spaces-resolve}
\externaldocument[formal-defos-]{formal-defos}
\externaldocument[defos-]{defos}
\externaldocument[cotangent-]{cotangent}
\externaldocument[examples-defos-]{examples-defos}
\externaldocument[algebraic-]{algebraic}
\externaldocument[examples-stacks-]{examples-stacks}
\externaldocument[stacks-sheaves-]{stacks-sheaves}
\externaldocument[criteria-]{criteria}
\externaldocument[artin-]{artin}
\externaldocument[quot-]{quot}
\externaldocument[stacks-properties-]{stacks-properties}
\externaldocument[stacks-morphisms-]{stacks-morphisms}
\externaldocument[stacks-limits-]{stacks-limits}
\externaldocument[stacks-cohomology-]{stacks-cohomology}
\externaldocument[stacks-perfect-]{stacks-perfect}
\externaldocument[stacks-introduction-]{stacks-introduction}
\externaldocument[stacks-more-morphisms-]{stacks-more-morphisms}
\externaldocument[stacks-geometry-]{stacks-geometry}
\externaldocument[moduli-]{moduli}
\externaldocument[moduli-curves-]{moduli-curves}
\externaldocument[examples-]{examples}
\externaldocument[exercises-]{exercises}
\externaldocument[guide-]{guide}
\externaldocument[desirables-]{desirables}
\externaldocument[coding-]{coding}
\externaldocument[obsolete-]{obsolete}
\externaldocument[fdl-]{fdl}
\externaldocument[index-]{index}

% Heap Project documents
\externaldocument[k3-]{k3}
\externaldocument[torelli-]{torelli}
\externaldocument[linear-algebra-]{linear-algebra}
\externaldocument[groups-]{groups}
\externaldocument[complex-analysis-]{complex-analysis}
\externaldocument[functional-analysis-]{functional-analysis}
\externaldocument[categories-]{categories}
\externaldocument[commutative-algebra-]{commutative-algebra}
\externaldocument[ringed-spaces-]{ringed-spaces}
\externaldocument[homological-algebra-]{homological-algebra}
\externaldocument[lie-algebras-]{lie-algebras}
\externaldocument[representation-theory-]{representation-theory}
\externaldocument[smooth-manifolds-]{smooth-manifolds}
\externaldocument[differential-topology-]{differential-topology}
\externaldocument[algebraic-topology-]{algebraic-topology}
\externaldocument[differential-geometry-]{differential-geometry}
\externaldocument[algebraic-geometry-]{algebraic-geometry}
\externaldocument[symplectic-geometry-]{symplectic-geometry}
\externaldocument[complex-geometry-]{complex-geometry}
\externaldocument[hodge-theory-]{hodge-theory}
\externaldocument[%
variations-of-hodge-structures-]{%
variations-of-hodge-structures}
\externaldocument[%
homological-mirror-symmetry-]{%
homological-mirror-symmetry}
\maybeexternaldocument{deformations-}{deformations}
\externaldocument[vertex-algebras-]{vertex-algebras}
\externaldocument[springer-]{springer}
\externaldocument[infty-categories-]{infty-categories}
\externaldocument[seminars-differential-geometry-]{%
seminars-differential-geometry}
\externaldocument[seminars-complex-geometry-]{%
seminars-complex-geometry}
\externaldocument[seminars-algebraic-geometry-]{%
seminars-algebraic-geometry}
\externaldocument[seminars-analysis-pde-]{%
seminars-analysis-pde}
\externaldocument[surfaces-]{surfaces}
\externaldocument[birational-maps-commutative-algebra-]{birational-maps-commutative-algebra}
\externaldocument[cremona-transformations-]{cremona-transformations}
\externaldocument[derived-categories-of-sheaves-]{derived-categories-of-sheaves}
\externaldocument[geometric-invariant-theory-]{geometric-invariant-theory}
\externaldocument[geometric-stability-conditions-and-group-actions-]{geometric-stability-conditions-and-group-actions}
\externaldocument[moduli-spaces-of-sheaves-]{moduli-spaces-of-sheaves}
\externaldocument[bridgeland-stability-general-theory-]{bridgeland-stability-general-theory}
\externaldocument[hilbert-schemes-]{hilbert-schemes}
\externaldocument[noncommutative-abelian-surfaces-and-kummer-type-hyperkahler-varieties-]{noncommutative-abelian-surfaces-and-kummer-type-hyperkahler-varieties}
\externaldocument[mumford-tate-groups-in-hodge-theory-]{mumford-tate-groups-in-hodge-theory}
\externaldocument[hodge-birational-atoms-2026-]{hodge-birational-atoms-2026}
\externaldocument[xxii-egd-2026-]{%
xxii-egd-2026}
\externaldocument[pde-]{pde}
\externaldocument[probability-]{probability}

\externaldocument[linguistics-]{linguistics}

\externaldocument[%
16th-alga-meeting-2026-]{%
16th-alga-meeting-2026}

\externaldocument[%
iv-jornadas-lefschetz-2026-]{%
iv-jornadas-lefschetz-2026}

\externaldocument[reading-the-masters-]{%
reading-the-masters}

\externaldocument[real-analysis-]{%
real-analysis}

\externaldocument[optimal-transport-]{%
optimal-transport}

\externaldocument[history-]{history}

\externaldocument[seiberg-witten-]{seiberg-witten}

\externaldocument[history-philosophy-biology-]{history-philosophy-biology}

\externaldocument[classical-mechanics-]{classical-mechanics}
\externaldocument[electromagnetism-]{electromagnetism}
\externaldocument[waves-and-optics-]{waves-and-optics}
\externaldocument[wif-]{waves-in-fluids-iv}
\externaldocument[thermal-physics-]{thermal-physics}
\externaldocument[quantum-mechanics-]{quantum-mechanics}
\externaldocument[special-relativity-]{special-relativity}
\externaldocument[physics-]{physics}

\externaldocument[chemistry-]{chemistry}

\externaldocument[molecular-cell-biology-]{molecular-cell-biology}

\externaldocument[microbiology-]{microbiology}

\externaldocument[protistology-]{protistology}

\externaldocument[mycology-]{mycology}

\externaldocument[botany-]{botany}

\externaldocument[zoology-]{zoology}

\externaldocument[animal-histology-]{animal-histology}

\externaldocument[animal-physiology-]{animal-physiology}

\externaldocument[systematics-biogeography-]{systematics-biogeography}

\externaldocument[ecology-natural-resources-]{ecology-natural-resources}

\externaldocument[earth-sciences-]{earth-sciences}

\externaldocument[biostatistics-]{biostatistics}

\externaldocument[mathematical-systems-biology-]{mathematical-systems-biology}

\externaldocument[cimpa-el-salvador-2024-]{cimpa-el-salvador-2024}

\externaldocument[anthropology-]{anthropology}

\externaldocument[yang-mills-theory-]{yang-mills-theory}

\externaldocument[seminars-algebra-]{seminars-algebra}

\externaldocument[seminars-physics-]{seminars-physics}

\externaldocument[seminars-computing-]{seminars-computing}

\externaldocument[oktobermat-2026-]{oktobermat-2026}

\externaldocument[saet-2026-]{saet-2026}

\externaldocument[seminars-symplectic-geometry-]{seminars-symplectic-geometry}

\externaldocument[seminars-dynamics-]{seminars-dynamics}

\externaldocument[oktobermat-2024-]{oktobermat-2024}

\externaldocument[quantum-groups-]{quantum-groups}

\externaldocument[legal-2027-]{legal-2027}

\externaldocument[derived-categories-]{derived-categories}

% Theorem environments.
%
\theoremstyle{plain}
\newtheorem{theorem}[subsection]{Theorem}
\newtheorem{proposition}[subsection]{Proposition}
\newtheorem{lemma}[subsection]{Lemma}
\newtheorem{corollary}[subsection]{Corollary}

\theoremstyle{definition}
\newtheorem{definition}[subsection]{Definition}
\newtheorem{example}[subsection]{Example}
\newtheorem{exercise}[subsection]{Exercise}
\newtheorem{situation}[subsection]{Situation}

\theoremstyle{remark}
\newtheorem{remark}[subsection]{Remark}
\newtheorem{remarks}[subsection]{Remarks}

\numberwithin{equation}{subsection}

% Macros
%
\def\lim{\mathop{\mathrm{lim}}\nolimits}
\def\colim{\mathop{\mathrm{colim}}\nolimits}
\def\Spec{\mathop{\mathrm{Spec}}}
\def\Hom{\mathop{\mathrm{Hom}}\nolimits}
\def\Ext{\mathop{\mathrm{Ext}}\nolimits}
\def\SheafHom{\mathop{\mathcal{H}\!\mathit{om}}\nolimits}
\def\SheafExt{\mathop{\mathcal{E}\!\mathit{xt}}\nolimits}
\def\Sch{\mathit{Sch}}
\def\Mor{\mathop{\mathrm{Mor}}\nolimits}
\def\Ob{\mathop{\mathrm{Ob}}\nolimits}
\def\Sh{\mathop{\mathit{Sh}}\nolimits}
\def\NL{\mathop{N\!L}\nolimits}
\def\CH{\mathop{\mathrm{CH}}\nolimits}
\def\proetale{{pro\text{-}\acute{e}tale}}
\def\etale{{\acute{e}tale}}
\def\QCoh{\mathit{QCoh}}
\def\Ker{\mathop{\mathrm{Ker}}}
\def\Im{\mathop{\mathrm{Im}}}
\def\Coker{\mathop{\mathrm{Coker}}}
\def\Coim{\mathop{\mathrm{Coim}}}

% Boxtimes
%
\DeclareMathSymbol{\boxtimes}{\mathbin}{AMSa}{"02}

%
% Macros for moduli stacks/spaces
%
\def\QCohstack{\mathcal{QC}\!\mathit{oh}}
\def\Cohstack{\mathcal{C}\!\mathit{oh}}
\def\Spacesstack{\mathcal{S}\!\mathit{paces}}
\def\Quotfunctor{\mathrm{Quot}}
\def\Hilbfunctor{\mathrm{Hilb}}
\def\Curvesstack{\mathcal{C}\!\mathit{urves}}
\def\Polarizedstack{\mathcal{P}\!\mathit{olarized}}
\def\Complexesstack{\mathcal{C}\!\mathit{omplexes}}
% \Pic is the operator that assigns to X its picard group, usage \Pic(X)
% \Picardstack_{X/B} denotes the Picard stack of X over B
% \Picardfunctor_{X/B} denotes the Picard functor of X over B
\def\Pic{\mathop{\mathrm{Pic}}\nolimits}
\def\Picardstack{\mathcal{P}\!\mathit{ic}}
\def\Picardfunctor{\mathrm{Pic}}
\def\Deformationcategory{\mathcal{D}\!\mathit{ef}}


\begin{document}

\title{Psychoanalysis}
\maketitle

\phantomsection
\label{section-phantom}

\tableofcontents

\section{Scope and sources}
\label{section-psychoanalysis-scope}

This notebook is a first introduction to
psychoanalysis. It goes from the big
picture to the details: what the word
means, where the ideas came from, the
main concepts of Freud, how a treatment
works, what came after Freud, and the
debate about its scientific status.

Freud is cited in the English Standard
Edition, edited by James Strachey. Its
translation choices (``instinct'' for
\emph{Trieb}, ``ego'' for \emph{Ich})
are themselves debated. For the meaning
of every technical term, the standard
reference is the dictionary of Laplanche
and Pontalis
\cite{laplanche-pontalis-language}. When
a term below is not given another
source, its entry there is the source.

These are study notes, not clinical
guidance. Claims about efficacy are
discussed separately, in
Section~\ref{section-psychoanalysis-science}.

\section{What psychoanalysis is}
\label{section-what-is-psychoanalysis}

Freud gave a three-part definition in
an encyclopedia article of 1922
\cite{freud-se-18-encyclopaedia}.
Psychoanalysis is the name of

\begin{enumerate}
\item a procedure for investigating
mental processes that are otherwise
almost inaccessible;
\item a method, based on that
investigation, for treating neurotic
disorders;
\item the body of psychological knowledge
obtained in this way.
\end{enumerate}

The three are tied together: the theory
is supposed to come out of what is
observed in treatment, and the treatment
is supposed to work by applying the
theory. This coupling is both the
strength of the discipline and the main
target of its critics (see
Section~\ref{%
section-psychoanalysis-science}).

The single idea behind everything is
that a large part of mental life is
unconscious, and that unconscious
wishes and conflicts have effects:
dreams, slips, symptoms and the
repetition of patterns in one's life.

\section{Origins: hysteria and the
talking cure}
\label{section-psychoanalysis-origins}

Freud was a neurologist in Vienna. In
the 1880s he studied hysteria, a
condition with bodily symptoms
(paralyses, loss of sensation, coughs)
without a detectable lesion.

With the physician Josef Breuer he
published \emph{Studies on Hysteria} in
1895 \cite{freud-se-2-hysteria}. Its
first case is Breuer's patient ``Anna
O.''. Her symptoms improved when, under
hypnosis, she talked about the
circumstances in which each one first
appeared. She called this the ``talking
cure''.

The book's thesis is that hysterical
patients suffer mainly from
reminiscences: memories charged with
emotion that were kept out of
consciousness. A symptom stands for such
a memory, and recalling it with its
emotion can relieve the symptom.

Freud then gave up hypnosis. He asked
patients instead to say everything that
came to mind, without selecting. This
is \emph{free association}, still the
basic rule of analysis. He also noticed
that patients resist remembering, and
named the force that keeps ideas out of
consciousness \emph{repression}.

\section{The unconscious}
\label{section-unconscious}

In psychoanalysis ``unconscious'' does
not mean only ``not attended to at the
moment''. Freud's first model of the
mind, the \emph{topographical model},
has three systems
\cite{freud-se-14-unconscious}.

\begin{enumerate}
\item The conscious: what one is aware
of now.
\item The preconscious: what is not
conscious now but can become so by
directing attention, such as a phone
number one knows.
\item The unconscious: contents kept
away from consciousness by repression,
which can reach it only in disguised
form.
\end{enumerate}

The unconscious, in this sense, follows
its own rules, which Freud called the
\emph{primary process}. It knows no
negation and no time, tolerates
contradiction, and moves emotional
charge freely from one idea to another.
Conscious thinking follows the
\emph{secondary process}: logic, delay,
reality.

\section{Dreams}
\label{section-dreams}

\emph{The Interpretation of Dreams}
(1900) is the founding book
\cite{freud-se-4-5-dreams}. Freud called
dreams the royal road to knowledge of
the unconscious. Its main claims are
these.

\begin{enumerate}
\item A dream is the fulfilment of a
wish (Chapter III).
\item One must distinguish the
\emph{manifest content}, the dream as
remembered, from the \emph{latent
content}, the thoughts behind it.
\item The \emph{dream-work} turns latent
thoughts into the manifest dream
(Chapter VI). Its main operations are
\emph{condensation}, where one image
stands for several thoughts, and
\emph{displacement}, where emotional
weight moves from an important idea to
a trivial one.
\item The disguise is the work of a
\emph{censorship}: a forbidden wish can
only appear distorted.
\end{enumerate}

Interpretation runs the dream-work
backwards. The dreamer associates freely
to each element of the dream, and the
latent thoughts are reconstructed from
these associations.

The same logic explains slips of the
tongue, forgetting of names and other
\emph{parapraxes} of daily life
\cite{freud-se-6-everyday}. A slip is a
compromise between an intention and a
disturbing, unconscious one.

\section{Symptoms as compromises}
\label{section-symptoms}

A neurotic symptom, like a dream or a
slip, is a \emph{compromise formation}.
It satisfies, in disguised form, both a
repressed wish and the force that
repressed it. This is why symptoms
persist even when they cause suffering:
they also bring a hidden gain.

So the method is not to suppress
symptoms directly, but to make their
meaning conscious. Freud's formula for
the aim of treatment is that what is
unconscious should become conscious.

\section{Drives and sexuality}
\label{section-drives}

Freud's \emph{Trieb} is translated
``drive'' (Strachey writes
``instinct''). A drive is a pressure
with a bodily source, an aim
(satisfaction) and an object. Unlike an
animal instinct, its object is not
fixed: it is the most variable part and
is found through one's history. The
energy of the sexual drives is called
\emph{libido}.

In the \emph{Three Essays on the Theory
of Sexuality} (1905)
\cite{freud-se-7-three-essays} Freud
claims that sexuality begins in infancy
and is broader than genital sexuality.
It starts as a set of partial drives
attached to different zones of the body.
In editions after 1905 and in later
papers he arranged these into stages of
development: oral, anal, phallic, a
latency period, and genital.

The \emph{Oedipus complex} is the
child's set of loving and hostile wishes
toward the parents, typically love for
one parent and rivalry with the other.
Freud first describes it in the
Interpretation of Dreams, in the section
on typical dreams of the death of
beloved persons
\cite{freud-se-4-5-dreams}. For him it is
the core of the neuroses. The way it is
resolved shapes later character and
relationships.

Late in his work Freud added a second
basic drive, aggressive and
self-destructive, the \emph{death
drive}, opposed to the life drives
(\emph{Eros}). This is one of the most
contested parts of his theory, and many
analysts do not accept it.

\section{Id, ego and superego}
\label{section-id-ego-superego}

In \emph{The Ego and the Id} (1923)
\cite{freud-se-19-ego-id} Freud replaced
the topographical model by the
\emph{structural model}. Its three
agencies are the following.

\begin{enumerate}
\item The \emph{id} (\emph{das Es},
``the it''): the reservoir of the
drives, entirely unconscious, ruled by
the pleasure principle.
\item The \emph{ego} (\emph{das Ich},
``the I''): the agency that perceives,
thinks, controls movement and mediates
between the id, the superego and
external reality. It is partly
unconscious.
\item The \emph{superego}
(\emph{das \"Uber-Ich}): the
internalized voice of the parents and
of society. It judges, forbids and sets
ideals. Guilt is the tension between
superego and ego. In Freud's account it
is the heir of the Oedipus complex
(Chapter III).
\end{enumerate}

The change of model was forced by a
clinical observation: the defences
themselves, and the sense of guilt, can
be unconscious. So ``conscious versus
unconscious'' could no longer be the
main division of the mind.

\section{Anxiety and defence}
\label{section-anxiety-defence}

In \emph{Inhibitions, Symptoms and
Anxiety} (1926)
\cite{freud-se-20-anxiety} Freud
reversed his earlier view. Anxiety is
not produced by repression; rather the
ego feels \emph{signal anxiety} when a
dangerous situation threatens, and this
signal sets off defences.

Anna Freud systematized the ego's
\emph{mechanisms of defence}
\cite{anna-freud-ego-defence}. The most
cited are listed below.

\begin{enumerate}
\item Repression: keeping an idea out
of consciousness.
\item Projection: attributing one's own
unacceptable wish to someone else.
\item Reaction formation: holding the
opposite attitude, such as excessive
kindness covering hostility.
\item Displacement: redirecting a
feeling to a safer object.
\item Isolation: keeping an idea but
cutting off its emotion.
\item Denial: refusing to acknowledge a
painful perception.
\item Sublimation: redirecting a drive
to a socially valued aim, such as art
or research.
\end{enumerate}

Defences are not pathological in
themselves. Everyone uses them; they
become a problem when they are rigid or
cost too much of one's life.

\section{The analytic situation}
\label{section-analytic-situation}

Classical analysis has a fixed
\emph{setting}: several sessions a week,
often four or five, of about fifty
minutes, often with the patient lying
on a couch and the analyst out of
sight. The fixed frame is meant to make
deviations from it meaningful. The main
notions of technique are these.

\begin{enumerate}
\item \emph{Free association}: the
patient says what comes to mind.
\item \emph{Evenly suspended
attention}: the analyst listens without
selecting in advance what is important.
\item \emph{Resistance}: everything that
opposes the work, such as silences,
forgetting and changing the subject. It
is not an obstacle to be overcome by
force, but material to be understood.
\item \emph{Transference}: the patient
repeats, with the analyst, feelings and
patterns from earlier relationships,
especially with the parents. Freud
first saw it as an obstacle, then as
the main tool of treatment, because the
past becomes present and can be
examined.
\item \emph{Countertransference}: the
analyst's own emotional reactions to
the patient. Freud saw it as something
to be mastered. Later analysts use it
as a source of information about the
patient.
\item \emph{Interpretation}: the analyst
puts into words a meaning the patient
has not yet recognized, often about the
transference.
\item \emph{Working through}: insight is
not enough. The same conflict has to be
recognized again and again in different
forms before anything changes.
\end{enumerate}

Today most treatments called
\emph{psychodynamic psychotherapy} use
these ideas with fewer sessions, face
to face.

\section{After Freud}
\label{section-after-freud}

Psychoanalysis split into schools. Each
keeps the unconscious and transference,
but puts something else at the center.

\emph{Ego psychology}, dominant in the
United States after the Second World
War, built on Anna Freud
\cite{anna-freud-ego-defence}. It
studies the ego's defences and its
adaptive functions.

Melanie Klein worked with young children
through play. She placed intense
phantasies and anxieties in the first
months of life. In her 1946 paper
\cite{klein-schizoid-1946} she describes
the \emph{paranoid-schizoid position},
in which the infant \emph{splits}
experience into wholly good and wholly
bad objects. It is followed by the
\emph{depressive position}, in which the
same person is recognized as both loved
and hated, with guilt and the wish to
repair. The same paper introduces the
idea of \emph{projective
identification}; the name itself was
added in its 1952 reprint.

\emph{Object relations} theory, and in
particular Donald Winnicott, moved the
focus from drives to relationships.
Winnicott's best-known idea is the
\emph{transitional object}
\cite{winnicott-transitional-1953}: the
blanket or soft toy that is the infant's
first possession that is neither fully
``me'' nor fully ``not-me''. It belongs
to an intermediate area of experience
which, for Winnicott, later becomes the
area of play, art and culture.

Jacques Lacan called for a ``return to
Freud'' read through linguistics. He
held that the unconscious is structured
like a language. His essay on the
\emph{mirror stage} (1949)
\cite{lacan-ecrits-fink} argues that the
infant forms its ``I'' by identifying
with its unified image in a mirror. So
the ego is built from an image of
oneself seen from outside. Lacanian
psychoanalysis is very influential in
France and Latin America, including
Brazil and Argentina.

\section{Is psychoanalysis a science?}
\label{section-psychoanalysis-science}

This question has been debated since
the beginning. Three positions are worth
knowing.

Karl Popper used psychoanalysis,
together with Marxism and Adler's
psychology, as his example of a theory
that is not falsifiable
\cite{popper-conjectures}. For him, any
behavior can be explained by it after
the fact, so no observation could
refute it. Its apparent confirmations
are therefore worthless.

Adolf Gr\"unbaum disagreed with Popper
\cite{grunbaum-foundations}. He argued
that Freud's theory does make testable
claims, but that the evidence gathered
on the couch cannot test them. The
analyst's suggestion can contaminate
the patient's associations and
agreement. So the clinical data do not
support the causal claims built on
them.

On treatment itself, the question is
empirical. Shedler's review of
meta-analyses
\cite{shedler-efficacy-2010} argues that
psychodynamic therapy has effect sizes
comparable to other evidence-based
therapies, and that benefits tend to
persist or grow after treatment ends.
This is a review by an advocate of the
method. The size and interpretation of
these effects remain debated, and the
efficacy of a therapy does not by
itself confirm the theory behind it.

A reasonable summary: several ideas
survived well outside psychoanalysis,
such as unconscious mental processes,
defences, and the influence of early
relationships. Specific claims, such as
the universality of the Oedipus complex
or the death drive, are not
established.

\section{How to continue}
\label{section-psychoanalysis-continue}

A short reading path, from the most
accessible to the most technical.

\begin{enumerate}
\item Freud's own \emph{Introductory
Lectures on Psycho-Analysis}
(1916--1917, Standard Edition volumes
15 and 16), written for a general
audience.
\item The chapters on dream-work in
\cite{freud-se-4-5-dreams}.
\item \cite{freud-se-19-ego-id}, which
is short.
\item The dictionary
\cite{laplanche-pontalis-language}, to
check any term precisely.
\item For the scientific debate,
\cite{grunbaum-foundations} and
\cite{shedler-efficacy-2010}.
\end{enumerate}

\bibliography{my}
\bibliographystyle{amsalpha}

\end{document}
